% BANKS-SOURCE: pdf=247 print=237 kind=prose
bottom boundaries of the square, but may be non-trivial on the right boundary.
The condition that $\phi$ be continuous is that

% BANKS-SOURCE: pdf=247 print=237 kind=display
\[
  g(1,t_2)\phi_0=\phi_0,
\]

for all $t_2$, i.e. $g(1,t_2)\in H$. This then defines a map of the circle into
$H$, or a member of $\Pi_1(H)$. It is easy to verify that the map we have just
constructed, namely that of $\Pi_2(G/H)$ into $\Pi_1(H)$, is a group
homomorphism.

The monopole and four-index epsilon formulas in this section are the $D=4$
Lorentzian specialization, with $\eta=(-,+,+,+)$ and orientation
$\epsilon^{0123}=+1$.

Conversely, given any element of $\Pi_1(H)$ that can be continuously distorted
into the identity when $\Pi_1(H)$ is embedded in $\Pi_1(G)$, we can construct
$g(t_1,t_2)$ in terms of the homotopy which distorts that element into the
identity. Finally, the image of the identity in $\Pi_2(G/H)$, the mapping
$g(t_1,t_2)=1$, is clearly the identity in $\Pi_1(H)$. Our map is a group
homomorphism, so this means that it is a one-to-one and onto mapping (an
isomorphism) between $\Pi_2(G/H)$ and the subgroup of $\Pi_1(H)$ which maps
into the identity in $\Pi_1(G)$. In particular, if $G$ is simply connected we
have an isomorphism between $\Pi_1(H)$ and $\Pi_2(G/H)$. The most interesting
case is when $\Pi_1(H)$ is an integral lattice, i.e. when $H$ is a product of
$U(1)$ factors. In this case, as in our $SO(3)/U(1)$ example, the topological
charge of solitons is precisely the magnetic charge of the multiple Maxwell
fields in the low-energy effective theory below the scale of the Higgs VEV. One
way to get such a pattern of symmetry breaking for any group is to introduce a
Higgs field in the adjoint representation. A general VEV for an adjoint Higgs
breaks the group to its Cartan torus.
\BanksImplicitHook{I-CH10-019}

Every Lie group $G$ has a simply connected covering group $\bar G$, such that
$G=\bar G/\Pi_1(G)$. Familiar examples are $\overline{SO(3)}=SU(2)$ and more
generally $\overline{SO(n,m)}=Sp(n,m)$. These examples illustrate a general
pattern: the covering group has extra representations, the spinors, with the
property that the eigenvalues of Cartan generators (the weights of the
representation) are fractions of the weights in bona-fide representations of
$G$. A more general class of examples is obtained by taking any simply
connected Lie group $G$ modulo its center $Z_G$. Only representations invariant
under the center are representations of $G/Z_G$.

The pure Yang--Mills Lagrangian contains only adjoint fields. Let us add Higgs
fields only in representations of $G/Z_G$. Then monopole charges are restricted.
The gauge group of the theory could be either $G$ or $G/Z_G$, but the Lagrangian
does not distinguish between them, so monopole charges are restricted to be in
the subgroup of $\Pi_1(H)$ which is trivial inside of $\Pi_1(G/Z_G)$. So, in our
$SO(3)/U(1)$ example, the monopole charges must all be even.

The physical meaning of all of this has to do with the Dirac quantization
condition. We will describe it explicitly only for the case of $SO(3)/U(1)$,
but the reader should be able to generalize it to all other cases. As we have
discussed in Chapter 8, the low-energy effective theory is Maxwell's
electrodynamics coupled to both electric and magnetic charges. This raises a
well-known problem, which was first solved by Dirac. An electrically charged
particle couples to the vector potential

% BANKS-SOURCE: pdf=247 print=237 kind=display
\[
  \int A_\mu\frac{\dd x^\mu}{\dd\tau},
\]

% BANKS-SOURCE: pdf=248 print=238 kind=prose
but the vector potential cannot be globally defined in the presence of magnetic
charge, because

% BANKS-SOURCE: pdf=248 print=238 kind=display
\[
  \epsilon^{\mu\nu\lambda\kappa}\partial_\nu F_{\lambda\kappa}
  =J_m^\mu.
\]

Dirac solved this problem for point-like monopoles by taking a particular
solution of this equation with $F$ non-zero only along a 2-surface running from
the world line of the monopole to infinity, and then introducing a vector
potential to describe the rest of the electromagnetic field. Essentially he
described a monopole as a semi-infinite solenoid. But we would like physics to be
independent of the choice of where we place this fictitious solenoid, if we want
the description of the interaction of monopoles and charges to be local.
Classically there is no problem, since the charged particle will never hit the
infinitely thin solenoid, but quantum mechanically we must make sure that there
is no Aharonov--Bohm phase when the charged particle follows a trajectory
encircling the solenoid. As we proved in Chapter 8, this leads to the Dirac
quantization condition

% BANKS-SOURCE: pdf=248 print=238 kind=equation id=10.23
\begin{equation}
  (e_i g_j-e_jg_i)=2\pi n_{ij}.
  \label{eq:10.23}
\end{equation}

We have written the condition for arbitrary pairs of particles, assuming that
each pair has both electric and magnetic charge. Each $n_{ij}$ must be an
integer.
\BanksImplicitHook{I-CH10-020}

Another illuminating way of deriving this condition is to compute the angular
momentum of the electromagnetic field of the particle pair. It's easy to verify
that this contains a term

% BANKS-SOURCE: pdf=248 print=238 kind=display
\[
  \delta L=\frac{e_i g_j-e_jg_i}{4\pi}\,\widehat{\mathbf r},
\]

where $\widehat{\mathbf r}$ is the unit vector pointing between the pair. Quantization of
angular momentum in half-integer units now implies that $n_{ij}$ is an integer.
We will see a striking consequence of this fact a bit later on.
\BanksImplicitHook{I-CH10-021}

For solitonic monopoles in theories in which $U(1)$ arises from spontaneous
breakdown of a simple group, both electric and magnetic charges are already
quantized. Monopole quantization follows from the topological conditions we have
been discussing, while electric charge quantization follows from the fact that
the weights of representations of compact Lie groups are discrete. The
restriction we discussed above, namely that the magnetic charge be in the
subgroup of $\Pi_1(H)$ which maps to the identity in $\Pi_1(G/Z_G)$, is the
statement that the monopole satisfies the Dirac quantization condition for
electric charges in all representations of $G$, including those which are not
representations of $G/Z_G$. For $SO(3)/U(1)$ the monopole charge is twice what
it needs to be to satisfy the Dirac quantization condition with particles in
tensor representations of $SO(3)$. If this were not so we would run into an
inconsistency when we tried to add fields in the spinor representations of
$SU(2)$ to the Lagrangian.

The discussion of charge quantization in the context of non-trivial classical
solutions of the field equations leads us to ask about electric charge carried by
classical field configurations. Why can't such classical charges take on
arbitrary continuous values? To answer this, we first choose a gauge, $A^a{}_{0}=0$.
In this gauge we canonically quantize

% BANKS-SOURCE: pdf=249 print=239 kind=prose
the spatial components of the gauge potential, and then impose the constraint
equation derived by varying $A^a{}_{0}$ as a constraint on physical states:

% BANKS-SOURCE: pdf=249 print=239 kind=display
\[
  D_i{}^{ab}\Pi_b{}^i+\epsilon^{abc}(\pi_b\phi_c-\pi_c\phi_b)
  \ket{\Psi_{\rm phys}}=0.
\]

$P_b{}^i=(1/g^2)\partial_tA_b{}^i$ is the canonical conjugate to the vector
potential and $\pi_a$ the canonical conjugate to the Higgs field. The operator
which must vanish on physical states is the generator of time-independent local
gauge transformations. The corresponding generator of global $U(1)$
transformations is

% BANKS-SOURCE: pdf=249 print=239 kind=display
\[
  Q=\int_{S^2}\hat\phi^a n_i\Pi_a{}^i.
\]

In this gauge, all charge-carrying configurations have time-dependent vector
potentials at infinity. Julia and Zee \cite{banks-ref-148} found the corresponding
configurations with both electric and magnetic charge in a different, static
gauge.

In $A_0=0$ gauge, the solutions have the form of a time-dependent gauge
transformation of the monopole solution. For small deviations, the behavior near
infinity is

% BANKS-SOURCE: pdf=249 print=239 kind=display
\[
  \delta A_i{}^a=\delta_i{}^a\frac{\theta(t)}{r^2},
\]

while the leading behavior of the Higgs field is unchanged. This is a normalizable
zero mode of the monopole solution and the collective coordinate $\theta(t)$
must be quantized. However, since it represents a $U(1)$-gauge rotation $\theta$
is a periodic variable. Its canonical momentum, the electric charge $Q$, is
therefore quantized. The minimal unit is the charge on the adjoint
representation. Quantization of $\theta$ thus leads to a discrete spectrum of
dyonic (or dual-charged) excitations of the monopole. The spacing between these
levels is of order $(\Delta Q)^2$ in units of the mass of the massive gauge
bosons.

A surprise is in store \cite{banks-ref-149,banks-ref-150} when we carry out the
collective coordinate quantization for monopoles in theories in which there are
fields in representations that transform under the center of $SU(2)$. The unit
of charge quantization is half what it was in the $SU(2)/\mathbb Z_2$ theory. As
noted above, if we calculate the angular momentum of some of the dyonic bound
states of this system, we will find half-integer values. Thus, we have
constructed spin-$\tfrac12$ particles from a purely bosonic system!

% BANKS-SOURCE: pdf=249 print=239 kind=section id=10.9
\subsection{Problems for Chapter 10}
